% ДРУК-01: автономний типографічний зразок, не матеріал для видавництва.
% Санскритські слова нижче є загальними лексемами, а не цитатою з видання.
\documentclass[12pt,a4paper]{article}
\usepackage{fontspec}
\usepackage{polyglossia}
\setdefaultlanguage{ukrainian}
\setmainfont{DejaVu Serif}
\newfontfamily\devanagarifont[Script=Devanagari]{Noto Serif Devanagari}
\usepackage[margin=26mm]{geometry}
\begin{document}
\section*{Шрифти української наукової книги}
\textbf{Статус:} технічна проба; не текст вірша і не перевірене видання.

\medskip
\noindent\textbf{Деванагарі (окремий шрифт):} {\devanagarifont आत्मा}\quad {\devanagarifont ज्ञानम्}\quad {\devanagarifont साक्षिन्}

\medskip
\noindent\textbf{IAST (перевірка діакритики):} ātmā; jñānam; sākṣin.

\medskip
\noindent\textbf{Українська читальна передача:} атман; джнянам; сакшін.

\medskip
\noindent\textbf{Науковий апарат:} тут має бути перевірене джерело, граматика, коментар і межі культурної аналогії. Відсутність джерела або дозволу не можна приховувати.
\end{document}
